\changecolours{ScoutGreen}
\section[Referências]{Referências Bibliográficas e Próximos Passos}

% ------------------------------------------------------------------------------
\begin{frame}{Referências Bibliográficas Fundamentais}
  \scriptsize
  \begin{itemize}\setlength{\itemsep}{1.5mm}
    \item \textbf{CHEN, Peter Pin-Shan.} \textit{The Entity-Relationship Model -- Toward a Unified View of Data}. ACM Transactions on Database Systems (TODS), v. 1, n. 1, p. 9--36, 1976.
    \item \textbf{ELMASRI, Ramez; NAVATHE, Shamkant B.} \textit{Sistemas de Banco de Dados}. 7. ed. São Paulo: Pearson Education do Brasil, 2018. (Capítulos 3 e 4: Modelagem Conceitual MER e EER).
    \item \textbf{HEUSER, Carlos Alberto.} \textit{Projeto de Banco de Dados}. 6. ed. Porto Alegre: Bookman, 2009.
    \item \textbf{SILBERSCHATZ, Abraham; KORTH, Henry F.; SUDARSHAN, S.} \textit{Sistema de Banco de Dados}. 7. ed. Rio de Janeiro: Elsevier, 2020.
  \end{itemize}
  \vspace{0.15cm}
  \alert{Próximo Encontro -- Aula 08 (20/10/2026 -- Terça-feira):}\par
  \scriptsize
  \textbf{Laboratório Prático de Modelagem Conceitual com Ferramentas CASE (\textit{brModelo} / \textit{Draw.io}):} Traga o seu caderno do Estudo Dirigido preenchido para construirmos o DER/EER oficial do seu projeto!
\end{frame}
